% Generated draft: requires mathematical and source review.
\documentclass[11pt]{article}
\usepackage[margin=25mm]{geometry}
\usepackage{fontspec}
\setmainfont{Latin Modern Roman}
\usepackage{amsmath,amssymb,mathtools,unicode-math}
\setmathfont{Latin Modern Math}
\usepackage{enumitem,needspace,xurl,hyperref}
\setlength{\parindent}{0pt}
\setlength{\parskip}{5pt}
\setlength{\emergencystretch}{3em}
\newcommand{\N}{\mathbb{N}}
\newcommand{\Z}{\mathbb{Z}}
\newcommand{\Q}{\mathbb{Q}}
\newcommand{\R}{\mathbb{R}}
\newcommand{\C}{\mathbb{C}}
\newcommand{\F}{\mathbb{F}}
\newcommand{\A}{\mathbb{A}}
\newcommand{\PP}{\mathbb{P}}
\newcommand{\E}{\mathbb{E}}
\newcommand{\eps}{\varepsilon}
\newcommand{\dd}{\,\mathrm{d}}
\begin{document}
\section*{UnsolvedMath ID: 701023}
\textbf{Source problem code:} catalogue entry
\subsection{Definitions and mathematical statement}
Let $p$ be an odd prime, let $\F_p$ be the field with $p$ elements, and equip the formal power series ring $\F_p[[t]]$ with its $t$-adic topology. Define
\[
u_p=\begin{pmatrix}1&t\\0&1\end{pmatrix},\qquad
v_p=\begin{pmatrix}1&0\\t&1\end{pmatrix},\qquad
H_p=\overline{\langle u_p,v_p\rangle}\leq \operatorname{SL}_2(\F_p[[t]]),
\]
where the bar denotes closure in the induced matrix topology. A pro-$p$ group is a topological group isomorphic to an inverse limit of finite $p$-groups. A group $G$ is soluble if its derived series $G^{(0)}=G$, $G^{(j+1)}=[G^{(j)},G^{(j)}]$ reaches the trivial group after finitely many steps. Write $\mathcal V_p(G)$ for the variety of pro-$p$ groups generated by $G$, using the convention that this is the smallest class of pro-$p$ groups containing $G$ and closed under closed subgroups, continuous quotient groups, and arbitrary direct products.

The supplied source attributes to A. N. Zubkov the following result for $2\times2$ matrices: for every odd prime $p$ and every pro-$p$ group $G$ of such matrices,
\[
G\text{ is soluble}\quad\text{or}\quad H_p\in\mathcal V_p(G).
\]
Here the matrix coefficient domain and its topology are not specified in the supplied record.

\textbf{Question (Kourovka Notebook 14.38).} Does the same conclusion hold in every matrix size $n\geq3$? More explicitly, for every odd prime $p$, every integer $n\geq3$, and every pro-$p$ group $G$ of $n\times n$ matrices in the class intended by the source, must
\[
G\text{ be soluble}\quad\text{or}\quad H_p\in\mathcal V_p(G)?
\]
The comparison group remains the explicitly displayed group $H_p\leq\operatorname{SL}_2(\F_p[[t]])$. The quantification over matrix groups cannot be made fully explicit from this record: the permitted coefficient fields or rings, the ambient matrix topology, and the requirements on the matrix realization must first be specified. The definition of $\mathcal V_p(G)$ above states the standard closure convention used in this draft; the source does not expand that terminology.
\subsection{Short English statement}
For an odd prime $p$, the source reports that a pro-$p$ group of $2\times2$ matrices is either soluble or generates a variety containing the specified two-generator closed subgroup of $\operatorname{SL}_2(\F_p[[t]])$. Zubkov asks whether this also holds for every matrix size at least $3$. The coefficient domain and topology of the matrix groups are unspecified. The registry labels the problem open; this is supplied status information, not an independently verified assessment.
\subsection{Sources}
\begin{itemize}
\item Supplied local UnsolvedMath registry record: permanent ID $701023$, exact source problem code \texttt{catalogue entry}, title \emph{Kourovka Notebook Problem 14.38}. The supplied status is \texttt{open}; the record attributes the proposal to A. N. Zubkov in 1999.
\item \emph{The Kourovka Notebook (Unsolved Problems in Group Theory)}, Problem 14.38, 14th issue (1999), as identified in the supplied record through the 21st issue edition (2026), \texttt{arXiv:1401.0300v48}. The supplied collection description names E. I. Khukhro and V. D. Mazurov as editors.
\item A. N. Zubkov, \emph{Algebra and Logic}, volume 29, no. 4 (1990), pp. 287--301. The supplied statement cites this reference for the $2\times2$ result. No article title is supplied, and neither the citation nor the precise hypotheses of that result have been independently verified here.
\end{itemize}
\end{document}
